\color{unidarkgreen}{\subsection{Максимальная токовая защита стороны ВН (МТЗ ВН)}}
\color{black}

\begin{enumerate}

\item
Максимальная токовая защита на стороне ВН с комбинированным пуском по напряжению предназначена для резервирования основных защит трансформатора и резервирования отключения КЗ на шинах среднего и низшего напряжения.

\item
Защита подключается к ТТ, установленным со стороны обмоток ВН силового трансформатора.

\item
Алгоритм работы МТЗ ВН приведен на рисунке \ref{pict:MTZ_VN}. Уставки МТЗ ВН приведены в таблице \ref{app_set:mtz}.

\begin{figure}[H]
\centering
\includegraphics[width=0.86\textwidth,height=0.86\textheight,keepaspectratio]{\fbpath/07.01 РЗТ/021.0311 МТЗ ВН/img/MTZ_VN.pdf}
\captionof{figure}{Функциональная схема логики МТЗ ВН}\label{pict:MTZ_VN}
\end{figure}


\item
В состав МТЗ ВН входят две ненаправленные ступени с контролем токов трех фаз высокой стороны трансформатора. 

\item
Порог срабатывания первой ступени защиты задается уставкой \textbf{«Iср МТЗ-1 ВН»}, а выдержка времени срабатывания первой ступени регулируется уставкой \textbf{«Тср МТЗ-1 ВН»}. Для повышения чувствительности, предусмотрена возможность пуска первой ступени МТЗ ВН от комбинированного пускового органа напряжения (КПОН) низкой и средней сторон трансформатора, которая задается программным переключателем \textbf{«Режим КПОН МТЗ-1 ВН»}. Ввод в работу первой ступени МТЗ ВН осуществляется с помощью программного переключателя \textbf{«МТЗ-1 ВН»}. 

\item
Порог срабатывания второй ступени защиты задается уставкой \textbf{«Iср МТЗ-2 ВН»}, а выдержка времени срабатывания второй ступени регулируется уставкой \textbf{«Тср МТЗ-2 ВН»}. Для повышения чувствительности, предусмотрена возможность пуска второй ступени МТЗ ВН от комбинированного пускового органа напряжения (КПОН) низкой и средней сторон трансформатора, которая задается программным переключателем \textbf{«Режим КПОН МТЗ-2 ВН»}. Ввод в работу второй ступени МТЗ ВН осуществляется с помощью программного переключателя \textbf{«МТЗ-2 ВН»}. 

\item
Оперативный ввод/вывод МТЗ ВН производится с помощью БУ (блока управления функциями) \textit{«Вывод МТЗ ВН»}. 

\item
Для силовых трансформаторов с заземленной или эффективно заземленной нейтралью во избежание неселективного действия при замыканиях на землю в сети 110-220 кВ рекомендуется измерительные органы МТЗ ВН включать на разность токов. С помощью программного переключателя \textbf{«ИО МТЗ ВН»} имеется возможность выбрать по фазным или линейным токам (по разности токов) будут работать измерительные органы первой и второй ступени МТЗ ВН. 

\item
Для обеспечения несрабатывания МТЗ ВН от бросков намагничивающего тока предусмотрена блокировка ступеней МТЗ по содержанию второй гармоники в фазном или линейном токе в зависимости от уставки \textbf{«ИО МТЗ ВН»}. Ввод в работу БНТ МТЗ ВН осуществляется с помощью уставки \textbf{«БНТ МТЗ ВН»}. Блокировка основана на контроле отношения действующего значения второй гармоники тока к основной гармонике. Ступени МТЗ блокируется при превышении данного соотношения пороговой величины, заданной уставкой \textbf{«I2г/I1г МТЗ ВН»}. Наличие функции блокировки от бросков намагничивающего тока задается для каждой ступени отдельно с помощью программных переключателей \textbf{«Режим БНТ МТЗ-1 ВН»} и \textbf{«Режим БНТ МТЗ-2 ВН»}. Для предотвращения излишней блокировки при больших токах КЗ, сопровождающихся насыщением ТТ, блокировка выводится из действия при превышении тока основной гармоники выше уставки \textbf{«Iср БНТ МТЗ ВН»}.

\item
Предусмотрена возможность оперативного ускорения МТЗ ВН. Срабатывание оперативного ускорения МТЗ ВН действует на отключение трансформатора с запретом АПВ. Ввод в работу ОУ МТЗ ВН и выбор ускоряемой ступени осуществляется с помощью программного переключателя \textbf{«ОУ МТЗ ВН»}, а выдержка времени срабатывания ОУ МТЗ ВН регулируется уставкой \textbf{«Тср ОУ МТЗ ВН»}. Оперативный ввод/вывод ОУ МТЗ ВН производится с помощью команды управления «Ввод ОУ МТЗ ВН».

\item
Действие срабатывания первой и второй ступени МТЗ ВН предусматривает отключение трансформатора со всех сторон с пуском УРОВ и запретом АПВ. 

\item
При опробовании силового трансформатора выключателем высокой стороны условия пуска МТЗ ВН от КПОН игнорируются. В этом случае пуск МТЗ ВН формируется по отсутствию сигналов включенного положения вводных выключателей низкой и средней стороны трансформатора.

\end{enumerate} 